\documentclass{beamer}

\usepackage{fontspec}
\usepackage[ngerman]{babel}
\usepackage{csquotes}
\usepackage{hyperref}

\title{Freie Software und Freies Wissen Dresden}
\subtitle{Fork us in real life!}
% \author{\url{http://fsfw-dresden.de/}}
\date{30.~Dezember 2017}

\usetheme{FSFW}

\begin{document}

\begin{frame}
  \maketitle
\end{frame}

\begin{frame}
  \frametitle{Wer sind wir?}
  \onslide<+->
  \begin{itemize}
  \item Hochschulgruppe an der TU Dresden (gegründet 2014, ca. 10 P.)
  \item Idee: \enquote{Lobbygruppe} für freie Software und freies Wissen
    \begin{itemize}
    \item Hochschulangehörigen als Zielgruppe (Multiplikationswirkung)
    \item Hochschule als Arbeitsfeld (Räume, Strukturen)
    \end{itemize}
    \bigskip\onslide<+->
  \item Bisherige Projekte
    \begin{itemize}
    \item Linux-Install-Party, Linux-Presentation-Day
    \item Verschlüsselungsgewinnspiel
    \item Crypto-Parties
    \item Monatliche Sprechstunde zu \LaTeX{} u.a.
    \item Formulierung eines Programmpapiers
    \item „Uni-Stick”:~100 $\times$ 8\,GB mit freier Software
    \item Git-Workshop
    \end{itemize}
  \end{itemize}
\end{frame}

\begin{frame}
  \frametitle{Warum Hochschulen als Arbeitsfeld?}
  \begin{itemize}
  \item<1-> Direkter Bezug zu freier Wissensverbreitung gegeben
  \item<2-> Problematische Monopolstellung einiger Wissenschaftsverlage
    ist evident
  \item<3-> Zum Teil findet \enquote{Produktschulung} statt
    \begin{itemize}
    \item[$\Rightarrow$] Abhängigkeit, \emph{vendor lock-in}
    \end{itemize}
  \item<4-> Hier werden zukünftige EntscheidungsträgerInnen ausgebildet
    und geprägt
  \end{itemize}
\end{frame}

\begin{frame}{Die Vision}
  \begin{itemize}
  \item<1-> Durch Schulung und Support die Verbreitung freier Software \\
    unter Studierenden erhöhen.
  \item<2-> Sensibilisierung für freie Software mit den Hochschulabsolventen  \\
    in die Wirtschaft tragen.
  \item<3-> Durch Projekte Reputation sammeln und diese an der Hochschule \\
    als politisches Kapital einsetzen.
  \item<4-> Breitenwirkung durch Forks an allen deutschen Hochschulen ;)
  \end{itemize}
\end{frame}

\begin{frame}{Fork us in real life!}
  \begin{itemize}
  \item<1-> zusammenfinden einer Gründergruppe (etwa drei Personen)
  \item<2-> Gründungstreffen planen, bewerben und abhalten
  \item<3-> regelmäßige Treffen organisieren
  \item<4-> Infrastruktur (Mailinglisten, Homepage, (a-)soziale Medien) aufsetzen
  \item<5-> erste Projekte angehen, z.\,B.
    \begin{itemize}
    \item Sprechstunde zu LaTeX, LibreOffice, etc.
    \item Crypto-Party
    \item Linux-Install-Party, Linux-Presentation-Day
    \item Positionspapier formulieren
    \end{itemize}
  \item<6-> \dots
  \item<7-> PROFIT!
  \end{itemize}
\end{frame}

\begin{frame}
  \frametitle{FSFW - Fork us in real life!}
  \begin{center}
    \Large Wir helfen gerne beim Forken!
  \end{center}
  \begin{columns}
    \column{0.4\textwidth}
    \includegraphics[width=\textwidth]{qr-website.pdf}
    \column{0.59\textwidth}
    \begin{description}
      \item[Web] \url{https://fsfw-dresden.de/}
      \item[Mail] \href{mailto:kontakt@fsfw-dresden.de}{\texttt{kontakt@fsfw-dresden.de}}
    \end{description}
  \end{columns}
\end{frame}

\end{document}
